\documentclass[lettersize,journal]{IEEEtran}
\usepackage{graphicx}
\usepackage[caption=false,font=normalsize,labelfont=sf,textfont=sf]{subfig}
\usepackage{textcomp}
\usepackage{verbatim}
\usepackage[comma,numbers,square,sort&compress]{natbib}
\usepackage{epstopdf}
\usepackage[cmex10]{amsmath}
\usepackage{array, arydshln, blkarray, enumerate, color, xcolor, multicol}
\usepackage{amsfonts, amsmath, amssymb, latexsym, algorithmic, algorithm, bm}
\usepackage{longtable, rotating, multirow, url}
\usepackage{booktabs}
\hyphenation{op-tical net-works semi-conduc-tor IEEE-Xplore}
\allowdisplaybreaks[4]
\newcommand{\field}[1]{\mathbb{#1}}
\newcommand{\calc}[1]{\mathcal{#1}}
\DeclareMathOperator{\PR}{\field{P}}
\DeclareMathOperator{\E}{\field{E}}
\DeclareMathOperator*{\argmax}{argmax}
\DeclareMathOperator*{\argmin}{argmin}
\def\N{\field{N}}
\def\R{\field{R}}
\def\A{\field{A}}
\def\X{\field{X}}
\def\W{\field{W}}
\def\Q{\field{Q}}
\def\I{\field{I}}
\def\K{\field{K}}
\def\V{\field{V}}
\def\cA{{\cal A}}
\def\cB{{\cal B}}
\def\cC{{\cal C}}
\def\cD{{\cal D}}
\def\cE{{\cal E}}
\def\cF{{\cal F}}
\def\cG{{\cal G}}
\def\cH{{\cal H}}
\def\cI{{\cal I}}
\def\cJ{{\cal J}}
\def\cK{{\cal K}}
\def\cL{{\cal L}}
\def\cM{{\cal M}}
\def\cN{{\cal N}}
\def\cO{{\cal O}}
\def\cP{{\cal P}}
\def\cQ{{\cal Q}}
\def\cS{{\cal S}}
\def\cT{{\cal T}}
\def\cU{{\cal U}}
\def\cV{{\cal V}}
\def\cW{{\cal W}}
\def\cX{{\cal X}}
\def\cY{{\cal Y}}
\def\cZ{{\cal Z}}

\usepackage{amsthm}
\newtheorem{theorem}{Theorem}
\newtheorem{definition}{Definition}
\newtheorem{lemma}{Lemma}
\newtheorem{proposition}{Proposition}
\newtheorem{corollary}{Corollary}
\newtheorem{remark}{Remark}

\begin{document}

\title{Feasibility-Aware Seed Budgeting for Constrained iLQR:\\
Diagnostic Multi-Armed Bandits and Bayesian Optimization}
\author{Jie Su}
\markboth{IEEE Robotics and Automation Letters,~Vol.~XX, No.~X, Month~2026}%
{J.~Su: Feasibility-Aware Seed Budgeting for Constrained iLQR}
\maketitle

\begin{abstract}
Constrained iterative LQR / DDP~\cite{tassa2012,howell2019,mastalli2020} is a workhorse for trajectory optimization in robotics, yet it is acutely sensitive to the initial trajectory: under a tight iteration budget, a poor warm-start yields slow convergence, local minima, or outright infeasibility. A standard remedy is multi-start: draw $K$ seeds, allocate each a uniform budget share, and retain the best feasible trajectory. Uniform allocation is simple---and wasteful---when only a minority of seeds lie in favourable basins. We formulate multi-start constrained iLQR as a sequential budget-allocation problem. After each short warm-started batch, the solver exposes diagnostics---feasible cost improvement $\Delta J$, constraint violation $v$, improvement rate $\rho$, and the closed-loop convergence measure $\chi$~\cite{zhu2024}---which form the selection signal: the primary reward combines the first three, while $\chi$ serves only as a termination switch and as a feature for Bayesian optimisation. We prove that any score in a natural class of local-improvement diagnostics can mis-rank seeds relative to their terminal feasible cost (ranking inversion), so greedy selection by a single score is not universally optimal. We instantiate Uniform, UCB, and GP--EI policies on this foundation. Adaptive allocation with the proposed reward improves over Uniform on a wheeled-navigation and an articulated-arm benchmark (feasible-first comparison on 5/5 random seeds; $\approx$38\% lower mean cost), while greedy $\chi$-selection proves unreliable; a contact-rich quadruped template provides a ranking-inversion diagnostic on a third physical domain. The framework is budget-exact; the benchmark definitions and all
allocation traces are released, and the scheduling framework and host
solver are distributed as a prebuilt binary for independent reproduction.
\end{abstract}

\begin{IEEEkeywords}
Trajectory optimization, constrained iLQR, warm-start, multi-armed bandit, Bayesian optimization, model predictive control.
\end{IEEEkeywords}

% ===========================================================================
\section{Introduction}
% ===========================================================================

\IEEEPARstart{M}{ulti-start} trajectory optimization is a common strategy for mitigating the
sensitivity of iterative LQR / DDP~\cite{tassa2012,howell2019,mastalli2020} to the initial guess.
Given a library of $K$ seed trajectories (drawn, e.g., via
LHS~\cite{mckay1979} or homotopy heuristics~\cite{merkt2021}),
a fixed total budget of $B$ solver iterations, and a uniform allocation
rule, each seed receives $B/K$ iterations and the best feasible result is
retained~\cite{gyorgy2011,jain2019}.

Uniform allocation is simple, but it is wasteful whenever the seeds differ
substantially in quality: a seed that starts in a poor basin receives the same
budget as one descending rapidly toward a good local optimum.  The decision
problem is therefore \emph{how to allocate the remaining budget among seeds
dynamically, based on observations collected during optimisation.}

The solver exposes several diagnostics after each short batch of $b$
iterations: the feasible cost $J$, the max constraint violation $v$, the
per-iteration improvement rate $\rho$, and the closed-loop convergence
measure $\chi = \|\prod \tilde{A}_k\|_2$~\cite{zhu2024}, where
$\tilde{A}_k = A_k - B_k K_k$ are the closed-loop dynamics under LQR
feedback.  These signals contain
information about the local landscape, but they are not a perfect predictor
of the terminal feasible cost $J^\star$ that a seed will achieve after full
budget.  We show formally (Section~\ref{sec:theory_ranking}) that a
\emph{ranking inversion} must occur: a score computed from batch diagnostics
can prefer a seed whose terminal feasible cost is worse than another that the
same score ranks lower.  This ranking inversion implies that greedy
selection by any fixed score is not universally optimal, and that some
form of exploration is needed in principle.

Against this backdrop we formulate multi-start constrained iLQR as a
sequential budget-allocation (multi-armed bandit) problem and instantiate
three policies---Uniform, UCB~\cite{auer2002}, and GP--EI (Bayesian
optimisation~\cite{srinivas2010,shahriari2016})---on a frozen
feasibility-improvement reward that combines $\Delta J$, $v$, and $\rho$,
with $\chi$ reserved as a termination switch and a feature for BO.

\textbf{Related work.} Constrained iLQR/DDP and public
solvers~\cite{tassa2012,howell2019,mastalli2020}; direct transcription
methods~\cite{posa2014,schulman2014}; multi-start and
warm-starts~\cite{gyorgy2011,jain2019,merkt2021}; MAB /
UCB~\cite{auer2002}; BO with GP--EI~\cite{shahriari2016,srinivas2010,
rasmussen2006}; the convergence measure $\chi$~\cite{zhu2024}.
Our contribution is a \emph{feasibility-aware budgeting} interface between
these literatures and the practical problem of allocating solver iterations
across warm-starts.

\textbf{Contributions.}
\begin{enumerate}
  \item \textbf{Ranking inversion analysis.} We define the class of
    local-improvement scores $S_{\mathrm{LI}}$ (Definition~\ref{def:li_scores})
    and prove that any score in this class can disagree with the
    oracle ordering by terminal feasible cost
    (Lemma~\ref{lem:li_incompleteness}).  This implies that greedy selection
    by such scores is not universally optimal
    (Corollary~\ref{cor:greedy_insufficiency}), motivating exploration.
  \item \textbf{Feasibility-aware reward and protocol.} We freeze a scalar
    primary reward (Equation~\eqref{eq:reward}) that combines feasible cost
    improvement, constraint violation, and improvement rate, with $\chi$
    serving only as a termination switch and a BO feature dimension.
    Uniform, UCB, and GP--EI policies share an exact budget-accounting
    scheduler with warm-start-only batches.
  \item \textbf{Empirical findings on real constrained-iLQR instances.}
    We demonstrate ranking inversion
    (Corollary~\ref{cor:greedy_insufficiency}) on wheeled-navigation and
    contact-rich quadruped benchmarks: both the feasible-improvement score
    and $-\log\chi$ mis-rank seeds against the oracle ordering
    (Tables~\ref{tab:ranking_flip} and~\ref{tab:quad_diag}).
    Under the frozen protocol, adaptive allocation (UCB, BO--EI,
    Greedy-$s_r$) wins the feasible-first comparison against Uniform on
    all five RNG seeds of the navigation benchmark (one-sided sign test,
    $p = 2^{-5}$) and attains $\approx$38\% lower mean cost on an
    all-feasible arm benchmark, while greedy $\chi$-selection is
    unreliable in both settings.
\end{enumerate}
\textbf{Scope.} The paper does not claim solver superiority: the host is a
constrained iLQR stack validated for parity with open-source solvers on a
standard gate suite (Section~\ref{sec:host_note}), and the scheduling
claims concern iteration-budget allocation. The quadruped template serves
as a ranking-inversion diagnostic, not as a policy-separation suite.
Adaptive policies (UCB, BO--EI, Greedy-$s_r$) may tie with each other; we
do not claim to separate them on the current seed libraries.
The implementation is C++17. The benchmark definitions, experimental
protocol, and all allocation traces are released; the seed-budgeting
scheduler and the constrained-iLQR host are distributed as a prebuilt
binary (Docker image) to enable independent reproduction.\footnote{\url{https://github.com/jiesu-zju/optimal-solver}}

% ===========================================================================
\section{Method: Feasibility-Aware Seed Budgeting}
% ===========================================================================

\subsection{Background: Constrained iLQR and Diagnostics}
\label{sec:chi_background}

Consider a discrete-time optimal control problem with horizon $N$, state dimension $n$, and control dimension $m$:
\begin{equation}
\label{eq:ocp}
\begin{aligned}
\min_{\{x_k, u_k\}} \; & \ell_N(x_N) + \sum_{k=0}^{N-1} \ell_k(x_k, u_k),\\
\text{s.t.}\quad & x_{k+1} = f_k(x_k, u_k), \\
                  & g_k(x_k, u_k) \leq 0, \quad \forall k \in [0, N-1],\\
                  & h_k(x_k, u_k) = 0, \quad \forall k \in [0, N-1], \\
                  & g_N(x_N) \leq 0, h_N(x_N) = 0.
\end{aligned}
\end{equation}
where $\ell_k(x_k, u_k)$ and $\ell_N(x_N)$ are the stage and terminal cost functions, $f_k(x_k, u_k)$ is the nonlinear dynamics, $g_k(x_k, u_k)$ and $g_N(x_N)$ are inequality constraints (control limits, joint limits, etc.), and $h_k(x_k, u_k)$ and $h_N(x_N)$ are equality constraints (waypoints, terminal pose, etc.). Constrained iLQR solves~\eqref{eq:ocp} via iterative linearization. At each iteration, the dynamics and cost are quadratized around the current nominal trajectory $\{\bar{x}_k, \bar{u}_k\}_{k=0}^N$, yielding a local LQR subproblem whose solution is obtained via backward Riccati recursion, producing feedback gains $K_k$ and feedforward terms $d_k$. Constraints are handled with an ADMM / Filter line-search stack and optional Box-QP control limits in the Riccati step (implementation details in the host solver; see the reproducibility footnote). The forward pass applies $\delta u_k = K_k \delta x_k + \alpha d_k$ with filter or backtracking line search.

The backward pass also yields the closed-loop convergence measure~\cite{zhu2024}:
\begin{equation}
\chi = \|\Phi\|_2, \quad \Phi = \prod_{k=N-1}^{0} \tilde{A}_k, \quad \tilde{A}_k = A_k - B_k K_k
\label{eq:chi_def}
\end{equation}
where $A_k = \partial f_k/\partial x$, $B_k = \partial f_k/\partial u$, and $K_k$ are the LQR gains.
The matrix $\tilde{A}_k$ is the closed-loop state transition matrix, and $\Phi$ propagates perturbations over the full horizon. $\chi < 1$ indicates exponential decay of trajectory perturbations (local basin of attraction), while $\chi \gg 1$ signals poor convergence. Computation requires one $O(N n^3)$ matrix chain multiplication---performed automatically at each backward pass with negligible overhead.  We return to the properties of $\chi$ in Section~\ref{sec:chi_analysis}.

\subsection{Seeds as Bandit Arms}

We cast multi-start iLQR as a multi-armed bandit~\cite{auer2002}: arm~$i$ is a
seed trajectory $\tau^{(i)}$. At round $t$ we select $i_t$, run a batch of $b$
iLQR iterations with \emph{trajectory} warm-start (a fresh solver instance;
dual variables are not assumed to persist), and observe diagnostics
$(J_{\mathrm{before}}, J_{\mathrm{after}}, v, \chi)$, where $v$ is the max
constraint violation over all stages.

\subsubsection{Primary reward design}

The reward encodes the decision objective: under a tight budget, prefer seeds that are feasible and whose feasible cost is still decreasing, while avoiding credit for cost reduction achieved under infeasibility.  Define

\begin{align*}
\Delta J &= J_{\mathrm{after}}-J_{\mathrm{before}},\\
\mathrm{viol}_+ &= \max(0,\,v-\varepsilon),\\
\Delta J_{\mathrm{feas}} &=
\begin{cases}
\Delta J, & v \le \varepsilon,\\
0, & \text{otherwise},
\end{cases}\\[2pt]
\rho &= \frac{\max(0,\,-\Delta J_{\mathrm{feas}})}{\max(1,\,b_{\mathrm{used}})}.
\end{align*}

The primary reward is

\begin{equation}
r^{(i)}_t = \mathrm{clip}\!\Big(-\Delta J_{\mathrm{feas}} - \lambda\,\mathrm{viol}_+ + \eta\,\rho\Big),
\label{eq:reward}
\end{equation}

with weights $\lambda,\eta>0$.  The three terms serve distinct roles:

\begin{itemize}
  \item $-\Delta J_{\mathrm{feas}}$ is the main signal, active only when the
    arm ends the batch feasible.  It rewards feasible cost reduction; an
    infeasible arm receives zero from this term.
  \item $-\lambda\,\mathrm{viol}_+$ penalises constraint violations
    linearly, remaining active whether or not the arm is feasible.
  \item $\eta\,\rho$ is a per-iteration rate bonus.  When the solver
    consumes fewer iterations than the batch size $b$ (early convergence or
    numerical failure), this term amplifies the improvement rate slightly.
\end{itemize}

\subsubsection{Role of $\chi$}

The convergence measure $\chi$ is \emph{not} inserted into~\eqref{eq:reward}.
Instead it serves two auxiliary purposes:
\begin{itemize}
  \item \textbf{Termination switch.} If $\chi>\chi_{\mathrm{cut}}$ or the batch
    fails numerically, arm~$i$ is suspended and removed from further allocation,
    preventing budget from being wasted on seeds whose local linearisation is
    unstable.
  \item \textbf{BO feature.} When Bayesian optimisation is used as the
    scheduling policy, $\log_{10}\chi$ is included as a dimension of the seed
    feature vector (Section~\ref{sec:bo_policy}), giving the GP surrogate
    access to convergence information without making $\chi$ the objective.
\end{itemize}

We also optionally stop the entire allocation once a feasible incumbent reaches
a declared near-optimal cost threshold (early success), avoiding expenditure of
residual budget on poor seeds.

\subsection{Theoretical Motivation: Why Instantaneous Arm Scores Are Not Fixed}
\label{sec:theory_ranking}

The scheduler observes a diagnostic score $s(d_i(t))$ after each warm-started batch and must decide which seed receives further iterations. Exploration (UCB / BO) is justified if greedy ranking by $s$ can disagree with the oracle ranking by terminal feasible cost $J_i^\star$. We establish this disagreement through a core lemma on local-improvement scores (Section~\ref{sec:li_incompleteness}), an auxiliary analysis of $\chi$ (Section~\ref{sec:chi_analysis}), and a decomposition of the primary reward $s_r$ that delimits its regime of applicability (Section~\ref{sec:reward_decomposition}).

\subsubsection{Local-improvement scores can mis-rank seeds}
\label{sec:li_incompleteness}

We first define the class of scores to which our main result applies.

\begin{definition}[Local-improvement score class]
\label{def:li_scores}
A score $s$ belongs to $S_{\mathrm{LI}}$ if it depends only on the diagnostics of a single warm-started batch and, for any two pulls where both arms finish feasible ($v \le \varepsilon$), $s$ is strictly increasing in the negative cost change $-\Delta J$ (larger feasible decrease $\Rightarrow$ higher score).
\end{definition}

The dominant term of our primary reward $s_r$---namely $-\Delta J_{\mathrm{feas}}$ when $v \le \varepsilon$---belongs to $S_{\mathrm{LI}}$.

\begin{lemma}[Local-improvement incompleteness]
\label{lem:li_incompleteness}
For any $s \in S_{\mathrm{LI}}$, there exist instances with two arms sharing the same initial feasible cost $J_0$ such that
\[
s(d_A(1)) > s(d_B(1)) \quad\text{but}\quad J_A^\star > J_B^\star,
\]
where $J_i^\star$ is the terminal feasible cost of arm~$i$ under equal remaining budget after round~1.
\end{lemma}

\begin{proof}
Fix $s \in S_{\mathrm{LI}}$, common initial cost $J_0$, and choose any $a > b > 0$ as the first-batch cost decreases of arms $A$ and $B$, respectively. Both batches finish feasible by construction, so $s(d_A(1)) > s(d_B(1))$ follows from $s \in S_{\mathrm{LI}}$.

Specify the second-batch decreases $c, d \ge 0$ as
\[
\begin{aligned}
J_A &: \; J_0 \xrightarrow{\;-a\;} J_0 - a \xrightarrow{\;-c\;} J_0 - a - c,\\
J_B &: \; J_0 \xrightarrow{\;-b\;} J_0 - b \xrightarrow{\;-d\;} J_0 - b - d.
\end{aligned}
\]
The terminal costs are $J_A^\star = J_0 - a - c$ and $J_B^\star = J_0 - b - d$.  The inequality $J_A^\star > J_B^\star$ is equivalent to $a + c < b + d$, i.e.\ $d - c > a - b > 0$.

Since $a - b$ is a fixed positive number, choose $c = 0$ (arm~$A$ makes no further progress after the first batch---a shallow basin) and any $d > a - b$ (arm~$B$ continues to improve---a deep basin).  These choices are always feasible, giving $s(d_A(1)) > s(d_B(1))$ with $J_A^\star > J_B^\star$. \hfill $\blacksquare$
\end{proof}

\begin{corollary}[Greedy insufficiency of single-batch local-improvement scores]
\label{cor:greedy_insufficiency}
No policy that greedily selects the arm with the highest current score $s \in S_{\mathrm{LI}}$ can be universally optimal with respect to terminal feasible cost $J^\star$.
\end{corollary}

\begin{remark}[Scope]
Corollary~\ref{cor:greedy_insufficiency} does not rule out the possibility that a richer score aggregating multi-batch history (e.g.\ a running average or a GP surrogate) could recover the oracle ordering.  It shows only that a single batch of local-improvement information is insufficient.  This distinction motivates the construction of multi-batch features in the BO policy (Section~\ref{sec:bo_policy}) and the exploration bonus in UCB---both of which accumulate information across pulls.
\end{remark}

\begin{remark}[Specialisation to $s_r$]
The dominant term of our primary reward $s_r$ (Equation~\eqref{eq:reward}) is $-\Delta J_{\mathrm{feas}}$.  When the batch finishes feasible ($v \le \varepsilon$), this term belongs to $S_{\mathrm{LI}}$ and therefore inherits the ranking inversion of Lemma~\ref{lem:li_incompleteness}.  The second term $-\lambda\,\mathrm{viol}_+$ is active only when the arm is infeasible; it is not covered by this lemma and is examined empirically in Section~\ref{sec:reward_decomposition}.
\end{remark}

\begin{remark}[Generality]
Lemma~\ref{lem:li_incompleteness} constructs an abstract pair of cost sequences; it does not guarantee that every constrained iLQR problem realises such a configuration.  Whether a concrete problem and seed library produce the required ``fast-into-shallow-basin'' and ``slow-into-deep-basin'' pair is an empirical question.  Scores outside $S_{\mathrm{LI}}$ (e.g.\ pure violation penalties) could in principle avoid the construction altogether.

The lemma suffices because $-\Delta J_{\mathrm{feas}} \in S_{\mathrm{LI}}$ is the dominant term driving our protocol's seed selection; the empirical evidence in Section~\ref{sec:ranking_flip} confirms that the required conditions arise in practice.
\end{remark}

\subsubsection{Note on $\chi$}
\label{sec:chi_analysis}

Recall $\chi = \|\prod \tilde{A}_k\|_2$ from Section~\ref{sec:chi_background}.
The spectral norm of a matrix product is not monotone in the norms of the
individual factors (e.g.\ $\operatorname{diag}(10,0.01)\cdot
\operatorname{diag}(0.01,10) = \operatorname{diag}(0.1,0.1)$ has norm
$0.1$, while tightening each factor to $\operatorname{diag}(9.9,0.01)$
yields a product with norm $98.01$).  Therefore $\chi$ is not a quantity
that behaves monotonically with local cost improvement---it can increase
when individual linearizations appear to contract.  In our protocol $\chi$
serves only as a termination switch (arms with diverging $\chi$ are
suspended) and as a BO feature dimension.

\subsubsection{Reward decomposition and regime of applicability}
\label{sec:reward_decomposition}

The primary reward $s_r = -\Delta J_{\mathrm{feas}} - \lambda\,\mathrm{viol}_+ + \eta\rho$ has three components whose roles differ in the ranking-inversion landscape.

\textbf{Feasible improvement ($-\Delta J_{\mathrm{feas}}$).}  This is the dominant term driving seed selection when an arm is feasible.  It belongs to $S_{\mathrm{LI}}$; Lemma~\ref{lem:li_incompleteness} therefore guarantees that ranking inversion can occur, and greedy selection by this term alone is not universally optimal.  Section~\ref{sec:ranking_flip} provides concrete illustrations of this inversion on real constrained-iLQR instances.

\textbf{Violation penalty ($-\lambda\,\mathrm{viol}_+$).}  This term does \emph{not} belong to $S_{\mathrm{LI}}$; it does not reward cost decrease and is active only in the infeasible regime.  It encodes a soft preference: arms that remain persistently infeasible accumulate a negative score that lowers their pull priority.  However, the term is not a provable separator between ``temporarily infeasible but promising'' and ``permanently infeasible'' arms, because a moderately infeasible arm with superficially good cost can still receive a higher total score than a barely feasible one.  Its effectiveness is problem-dependent and examined empirically in the experiments.

\textbf{Rate bonus ($\eta\rho$).}  This term provides a small per-iteration improvement bonus that becomes relevant when the solver converges early or fails numerically (so $b_{\mathrm{used}} < b$).  Under nominal operation its effect on ranking is minor.

\textbf{Regime of applicability.}  Lemma~\ref{lem:li_incompleteness} establishes that local-improvement scores---including the dominant term of $s_r$---are \emph{not} universally reliable predictors of terminal feasible cost.  This negative result does not, however, imply that $s_r$ is useless as a bandit reward.  Bandit policies (UCB, BO) require only that the reward distribution carries \emph{some} signal about arm quality.

In practice, the regimes where $s_r$ provides a useful signal can be characterised per seed library.  After warm-up, if the empirical mean rewards $\bar{s}_r^{(i)}$ (or the feasible incumbent costs) are positively correlated with the oracle ordering $J_i^\star$, then bandit allocation can exploit this signal.  When the correlation is weak---for example because warm-up already identifies the best basin, as on the current benchmarks---the adaptive policies remain tied \emph{with each other}, and harder seed libraries are required to separate them; the separation of the adaptive bloc from Uniform on the navigation and arm benchmarks (Section~\ref{sec:budgeting_results}) shows that $s_r$ nevertheless carries exploitable signal in practice.  When the correlation is systematically negative (early $s_r$ anti-correlates with $J^\star$ and later batches do not correct it), the reward or feature representation must be revised.  This per-library verification is part of the experimental protocol rather than a theoretical guarantee; Section~\ref{sec:budgeting_results} provides diagnostic statistics.

\subsection{UCB Policy (Independent Arms)}

UCB treats each seed as an independent bandit arm under reward~\eqref{eq:reward}. At scheduling round $t$, the policy selects the seed maximizing:
\begin{equation}
\begin{aligned}
i_t^{\text{UCB}} &= \argmax_{i \in \mathcal{A}_t} \Bigl( \hat{r}_{i,t}
   + c \sqrt{\frac{2\log t}{n_{i,t}}} \Bigr)
\end{aligned}
\label{eq:ucb_chi}
\end{equation}
where $\mathcal{A}_t \subseteq \{1,\dots,K\}$ is the set of active (non-terminated) seeds, $\hat{r}_{i,t}$ is the empirical mean reward, $n_{i,t}$ is its pull count, and $c$ is an exploration constant. Seeds with $\chi > \chi_{\text{term}}$ are terminated and removed from $\mathcal{A}_t$. The total budget $B$ is the only hyperparameter.

\noindent\textbf{Theoretical justification.}
With clipped feasible-improvement rewards and $\chi$-based arm suspension, active-arm rewards are bounded. We use $c=1.0$.

\subsection{BO Policy (GP--EI)}
\label{sec:bo_policy}

While UCB treats seeds independently, BO models similarity via a Gaussian Process and selects by Expected Improvement~\cite{srinivas2010} on reward~\eqref{eq:reward}.  We use it as a tool, not a methodological innovation; the description here is kept brief.

Each arm is mapped to a diagnostic feature vector
$\phi^{(i)} = [\, \Delta J,\; v,\; \log_{10}\chi,\; n_i,\; \bar\rho \,]$,
updated after every pull.  The arm index is excluded to prevent the GP from
learning arbitrary arm identity rather than generalisable features.

We model the reward function over seed feature space as a GP~\cite{rasmussen2006} with RBF kernel $k(\phi, \phi') = \sigma_f^2 \exp(-\|\phi - \phi'\|_2^2 / 2\ell^2)$,
kernel hyperparameters set to the benchmark values listed in Table~\ref{tab:solver_params}---treated as fixed settings rather than optimised values.  The GP posterior mean $\mu(\phi^*)$ and variance $\sigma^2(\phi^*)$ follow the standard closed-form GP equations~\cite{rasmussen2006}; exact Cholesky inference costs $O(K^3)$, negligible for $K=8$.

The next seed is selected by maximising Expected Improvement~\cite{srinivas2010}
$\text{EI}(\phi) = (\mu - r_{\text{best}} - \xi)\,\Phi(Z) + \sigma\,\varphi(Z)$
with $Z = (\mu - r_{\text{best}} - \xi)/\sigma$ and $\xi = 0.01$, where $\Phi$ and $\varphi$ are the standard normal CDF and PDF.

\begin{algorithm}[t]
\caption{Feasibility-Aware Seed Budgeting}
\label{alg:budgeting}
\begin{algorithmic}[1]
\STATE \textbf{Input:} Seeds $\{\tau^{(i)}\}_{i=1}^K$, budget $B$, batch $b$, warm-up $w$, policy $\pi \in \{\text{UCB}, \text{BO}\}$
\STATE Initialize $n_i \gets 0$, $\hat{r}_i \gets 0$, $\chi_i \gets \infty$ for all $i$
\STATE \textbf{Phase 1:} Run $w$ warm-started iLQR iterations on each seed; compute diagnostics and $r_i$ via~\eqref{eq:reward}; update $n_i \gets n_i + 1$
\FOR{$t = Kw+1$ \textbf{to} $B$ \textbf{step} $b$}
    \STATE $i^* \gets \textsc{Select}(\pi)$ \quad // UCB or EI policy
    \STATE Run $b$ warm-started iLQR iterations on seed $i^*$ (using the previous solution as the initial guess); observe $\chi_{i^*}$
    \STATE Observe diagnostics; set $r_{i^*}$ via~\eqref{eq:reward}; update $\hat{r}_{i^*}$, $n_{i^*}$
    \IF{$\chi_{i^*} > \chi_{\text{term}}$} \STATE Terminate seed $i^*$ \ENDIF
\ENDFOR
\STATE \textbf{Return} solution of seed with lowest final cost
\end{algorithmic}
\end{algorithm}

% ===========================================================================
% ===========================================================================
\section{Experiments}
\label{sec:experiments}

We evaluate feasibility-aware seed budgeting on three frozen instances that
collectively cover the paper's theoretical claims: ranking inversion
(Layer~1), policy discrimination under $s_r$, and the negative control
$\chi$-greedy.  All experiments use reward~\eqref{eq:reward} with
$(\lambda,\eta)=(10,0.1)$, $\chi_{\mathrm{term}}=100$, disabled
early-success stopping, and the solver parameters of
Table~\ref{tab:solver_params}.

% ===========================================================================
\subsection{Setup and Protocol}
\label{sec:setup}
% ===========================================================================

Table~\ref{tab:protocol} lists the three benchmarks.  Each is run with
five random seeds ($42$--$46$), and reported metrics are mean $\pm$ std
across those seeds unless noted.

\begin{table*}[t]
\caption{Benchmark protocol.}
\label{tab:protocol}
\centering
\footnotesize
\setlength{\tabcolsep}{3pt}
\begin{tabular}{lcccccl}
\toprule
ID & System & $K$ & Warm-up & Batch & $B$ & Role \\
\midrule
B1 & Nav2D--Hard (bicycle + dense soft obstacles) & 8 & 10 & 10 & 160 & Primary policy discrimination \\
B2 & Planar quadruped (soft-contact template)$\dagger$ & 8 & 2 & 5 & 80 & Contact-rich ranking-flip diagnostic \\
B3 & 5-DOF planar analytic arm (UP/DOWN homotopy) & 8 & 3 & 5 & 80 & Second discrimination + clean $\chi$ negative control \\
\bottomrule
\end{tabular}

{\footnotesize $\dagger$ Diagnostic mode uses $K{=}4$ seeds with equal
per-seed budget (warm-up 2, then $+30$ iterations) and no bandit; the
$K{=}8$ suite configuration is a near-tie sanity check only.}
\end{table*}

\noindent\textbf{Policies.}
We compare five strategies sharing reward~\eqref{eq:reward}:
\textbf{Uniform} (equal allocation), \textbf{UCB}~\cite{auer2002}
(Eq.~\eqref{eq:ucb_chi}, $c=1.0$), \textbf{BO--EI} (GP with RBF kernel,
Section~\ref{sec:bo_policy}), \textbf{Greedy-$s_r$} (always pull the seed
with the highest current $s_r$), and \textbf{Greedy-$\chi$} (always pull the
seed with the lowest $\chi$).  The scheduler follows
Algorithm~\ref{alg:budgeting} with trajectory warm-start and $\chi_{\mathrm{term}}=100$.

\subsubsection{Benchmark descriptions}

\textbf{B1: Nav2D--Hard.}
Dynamics: bicycle model~\cite{zhu2024} with state
$x = [x,y,\theta,\kappa]^\top$ and control $u = [v,\dot\kappa]^\top$,
discretised via RK4 ($\Delta t = 0.1$~s, $N = 80$).
Cost: waypoint attraction $w_p\|p - p^*\|^2$, curvature penalty
$w_\kappa\kappa^2$, control regularisation $w_u\|u\|^2$, plus six circular
soft obstacles with quadratic penetration cost
$W\max(0,\,r+m-\|p-c\|)^2$ ($W=80$, $m=0.35$~m) forming a dense field with
two homotopy classes.
Constraints: control limits $|v| \le 3$, $|\dot\kappa| \le 2$; curvature
$|\kappa| \le 1.5$; terminal position tolerance $\pm 0.5$~m, heading
$\pm 0.3$~rad.
Setup: start $(0,0)$, target $(10,0)$; $K=8$ seeds from LHS noise
($w=0.5$, alternating heading bias), $B=160$, $b=10$, warm-up $10$.

\textbf{B2: Planar quadruped.}
Dynamics: planar floating-body template with state
$(p_x,p_z,\theta,q_{\mathrm{hip}}^\mathrm{f},q_{\mathrm{knee}}^\mathrm{f},
q_{\mathrm{hip}}^\mathrm{b},q_{\mathrm{knee}}^\mathrm{b},\dot q) \in \R^{14}$
and control (joint torques) $u \in \R^4$; analytic Jacobians.
Cost: quadratic body-pose tracking and joint regularisation; ground
interaction and joint limits encoded as soft costs (step penetration
penalty, hip retraction barrier, knee-angle bounds) rather than explicit
contact forces.
Constraints: knee-angle bounds; terminal box on body position
$|p_x| \le 0.2$, $-0.1 \le p_z \le 0.4$.
Setup: $K=8$ seeds, $B=80$, $b=5$, warm-up $2$; for the ranking-flip
diagnostic (no bandit), $K=4$, warm-up $2$, $+30$ additional iterations.

\textbf{B3: Planar5.}
Dynamics: 5-DOF planar analytic arm, double-integrator discretisation,
state $x \in \R^{10}$, control $u \in \R^5$, exact Jacobians.
Cost: quadratic end-effector tracking, joint regularisation, soft disk
obstacle.
Constraints: joint-limit inequalities; no terminal box.
Setup: $K=8$ seeds split into an \textit{UP} cluster (bait: fast initial
descent into a shallow basin) and a \textit{DOWN} cluster (gem: slower
initial descent into a deeper terminal basin), directly instantiating
Lemma~\ref{lem:li_incompleteness}; $B=80$, $b=5$, warm-up $3$.

\begin{figure}[t]
\centering
\includegraphics[width=\columnwidth]{figures/fig_bench_scenes.pdf}
\caption{Benchmark schematics (robot geometry first). Faint dotted mists mark
\emph{barrier-cost ridges} in~$J$ (formulation detail), not obstacle bodies.
(a)~Nav2D--Hard: left/right homotopy warm-starts around a mid-field cost ridge.
(b)~Planar quadruped with torso and thigh/shank capsules; seed~A over-lifts
the front foot (bait) while seed~B keeps a mild clearance above the ground
clearance-cost zone (gem).
(c)~Planar5 solid-link arm in a left-side ready pose; UP (bait) / DOWN (gem)
end-effector corridors flank a mid-workspace cost ridge.}
\label{fig:bench_scenes}
\end{figure}

\begin{table}[t]
\caption{Solver and scheduling parameters.}
\label{tab:solver_params}
\centering
\footnotesize
\setlength{\tabcolsep}{3pt}
\begin{tabular}{lcc}
\toprule
\textbf{Parameter} & \textbf{Symbol} & \textbf{Value} \\
\midrule
Reward: feasible-improvement coefficient & $1$ & (implicit) \\
Reward: violation weight & $\lambda$ & 10 \\
Reward: rate-bonus weight & $\eta$ & 0.1 \\
UCB exploration constant & $c$ & 1.0 \\
$\chi$ termination threshold & $\chi_{\mathrm{term}}$ & 100 (never triggered) \\
Early-success cost threshold & $J_{\mathrm{succ}}$ & $\infty$ (disabled) \\
GP kernel length scale & $\ell$ & 2.0 \\
GP signal variance & $\sigma_f^2$ & 1.0 \\
GP observation noise & $\sigma_n^2$ & $10^{-4}$ \\
El hyperparameter & $\xi$ & 0.01 \\
\midrule
ADMM inequality penalty & $\mu$ & 10 \\
Regularization & $\lambda_{\mathrm{reg}}$ & $10^{-4}$ \\
Constraint tolerance & $\varepsilon$ & $10^{-3}$ \\
Line search reduction & $\alpha_{\mathrm{decay}}$ & 0.5 \\
Filter margin & $\delta$ & $10^{-4}$ \\
Filter max size & $M$ & 5 \\
\bottomrule
\end{tabular}
\end{table}

% ===========================================================================
\subsection{Host Solver}
\label{sec:host_note}
% ===========================================================================

All experiments use the same constrained iLQR stack described in
Section~\ref{sec:chi_background}.  Because the host is distributed as a
prebuilt binary rather than source, we validate its behaviour on a
standard gate suite of five problems (pendulum swing-up, double-integrator
reach, cartpole stabilisation, quadrotor hover, bicycle navigation) with
three seeds each.  All 15 instances converge to feasible solutions, and
the resulting costs match the best open-source solver among
Crocoddyl~\cite{mastalli2020}, Aligator, and ALTRO~\cite{howell2019} to
within 2\% (Table~\ref{tab:gate}); the baseline scripts for the
open-source solvers are part of the release.  This establishes the host as
a credible reference implementation; the paper's claims concern
scheduling, not solver superiority.

\begin{table}[t]
\caption{Gate validation of the host against the best feasible
open-source baseline (Crocoddyl, Aligator, or ALTRO); five problems
$\times$ three seeds, all 15/15 instances feasible.}
\label{tab:gate}
\centering
\footnotesize
\begin{tabular}{lcc}
\toprule
Problem & Ours & Best SOTA (ratio) \\
\midrule
Pendulum swing-up & 186.1 & 1.000 \\
Double-integrator reach & 6.67 & 1.000 \\
Cartpole stabilise & 35--114 & $\approx$1.000 \\
Quadrotor hover & 4431--5477 & $\le$1.02 \\
Bicycle navigation & 538.9--539.5 & $\approx$1.000 \\
\bottomrule
\end{tabular}
\end{table}

% ===========================================================================
\subsection{Ranking-Flip Diagnostics}
\label{sec:ranking_flip}
% ===========================================================================

Lemma~\ref{lem:li_incompleteness} establishes that, at the level of
abstract cost sequences, local-improvement scores in $S_{\mathrm{LI}}$
can mis-rank seeds.  We verify that this inversion occurs on real
constrained-iLQR problems.

\subsubsection{Nav~2D equal-budget ranking flip}

On the Gate-aligned bicycle navigation benchmark with equal per-arm budgets
(80 iterations each, no bandit), the first batch of 10 iterations already
mis-ranks the seeds (Table~\ref{tab:ranking_flip}): greedy selection by
$s_\chi$ picks seed~1 and greedy selection by $s_r$ picks seed~0, whereas
the oracle ordering by terminal feasible cost prefers seed~2
($J^\star = 538.88$).  Both scores therefore invert against the oracle.

\begin{table}[t]
\caption{Ranking inversion on Nav homotopy seeds.  Equal per-arm budget
(80 iterations each), no bandit.  Scores are those observed after the first
10-iteration batch: $s_\chi = -\ln\chi$; $s_r$ is the clipped
feasible-improvement reward, with ties at the clip cap (50) broken by
first-seen order as in the scheduler.}
\label{tab:ranking_flip}
\centering
\footnotesize
\begin{tabular}{lrrrcc}
\toprule
Seed & Oracle $J^\star$ & $s_\chi$ & $s_r$ & Greedy-$\chi$ & Greedy-$s_r$ \\
\midrule
0 & 871.26 & 4.03 & 50.0 & & $\checkmark$ \\
1 & 538.91 & 6.33 & 50.0 & $\checkmark$ & \\
2 & \textbf{538.88} & 5.83 & 6.35 & & \\
\bottomrule
\end{tabular}
\end{table}

\subsubsection{Contact-rich ranking flip (planar quadruped)}

The same early-vs-late inversion appears on the quadruped template
(Table~\ref{tab:quad_diag}).  Cluster~A (over-lifting gait) looks best
after a short warm-up; cluster~B (mild clearance) attains a lower terminal
cost once more iterations are granted.  The gate passes on all five RNG
seeds.

\begin{table*}[t]
\caption{Contact-rich ranking-flip diagnostic (planar quadruped template,
no bandit).  Each of $K=4$ seeds receives warm-up 2 then $+30$ iLQR iterations.
Cluster A (over-lift) is better after warm-up; cluster B (mild clearance)
achieves lower terminal cost.}
\label{tab:quad_diag}
\centering
\footnotesize
\begin{tabular}{lrrrl}
\toprule
RNG & Best A @warm & Best A @late & Best B @late & Gate \\
\midrule
42 & 22.2 & 22.2 & \textbf{7.98} & PASS \\
43 & 22.9 & 22.9 & \textbf{9.39} & PASS \\
44 & 26.0 & 26.0 & \textbf{7.95} & PASS \\
45 & 23.1 & 23.1 & \textbf{7.98} & PASS \\
46 & 22.1 & 22.1 & \textbf{13.1} & PASS \\
\bottomrule
\end{tabular}
\end{table*}

These diagnostics instantiate the local-improvement incompleteness argument
on three physically distinct systems (wheeled, legged, and articulated);
they do not by themselves prove that any particular bandit policy dominates
Uniform.

% ===========================================================================
\subsection{Seed-Budgeting Results}
\label{sec:budgeting_results}
% ===========================================================================

\subsubsection{B1: Nav2D--Hard}

Under the frozen protocol, adaptive allocation with $s_r$ improves over
Uniform on the dense-obstacle navigation benchmark
(Table~\ref{tab:b1_results}).  Because many terminal iterates remain slightly
infeasible at $\varepsilon=10^{-3}$, we rank policies by feasible-first order
(feasible $\succ$ lower violation $\succ$ lower cost).  The adaptive bloc
(UCB, BO--EI, Greedy-$s_r$) wins the feasible-first ranking against Uniform
on all five RNG seeds; under the null hypothesis that the per-seed outcome
is decided by chance, five wins in five seeds has probability $2^{-5}$
(one-sided sign test, $p \approx 0.03$).  Greedy-$\chi$ attains a lower mean
cost but at the cost of feasibility---it acts as a negative control that
confirms the $\chi$ signal is not a reliable primary objective.

\begin{table}[t]
\caption{Nav2D--Hard seed budgeting ($K=8$, $B=160$, $n=5$).  Adaptive
policies (UCB, BO--EI, Greedy-$s_r$) tie with each other and win the
feasible-first ranking against Uniform on 5/5 seeds.}
\label{tab:b1_results}
\centering
\footnotesize
\setlength{\tabcolsep}{4pt}
\begin{tabular}{lccc}
\toprule
Policy & Cost (mean $\pm$ std) & Feas. rate & Viol. (mean) \\
\midrule
Uniform & $2002.6 \pm 435.6$ & 0/5 & 0.306 \\
UCB & $1587.5 \pm 852.4$ & 2/5 & 0.083 \\
BO--EI & $1587.5 \pm 852.4$ & 2/5 & 0.083 \\
Greedy-$s_r$ & $1587.5 \pm 852.4$ & 2/5 & 0.083 \\
Greedy-$\chi$ & $1418.7 \pm 314.4$ & 1/5 & 0.234 \\
\bottomrule
\end{tabular}
\end{table}

\subsubsection{B3: Planar5}

On the fully-feasible planar arm benchmark (Table~\ref{tab:b3_results}),
every run terminates without constraint violation, isolating score
mis-ranking from feasibility entanglement.  BO--EI and Greedy-$s_r$
reach mean cost $269.7 \pm 143.4$, approximately $38\%$ lower than
Uniform ($433.9 \pm 89.9$); paired over the five RNG seeds, every adaptive
policy attains lower cost than Uniform on every seed (one-sided sign test,
$p = 2^{-5}$).  Greedy-$\chi$ remains worse at
$402.4 \pm 50.5$ despite full feasibility---a clean negative control.

\begin{table}[t]
\caption{Planar5 seed budgeting ($K=8$, $B=80$, $n=5$, all runs feasible).
BO--EI and Greedy-$s_r$ recover the DOWN homotopy; Uniform often stays on
the UP bait; Greedy-$\chi$ is feasible but worse.}
\label{tab:b3_results}
\centering
\footnotesize
\setlength{\tabcolsep}{4pt}
\begin{tabular}{lcc}
\toprule
Policy & Cost (mean $\pm$ std) & Feas. rate \\
\midrule
Uniform & $433.9 \pm 89.9$ & 5/5 \\
UCB & $277.0 \pm 138.4$ & 5/5 \\
BO--EI & $\mathbf{269.7 \pm 143.4}$ & 5/5 \\
Greedy-$s_r$ & $\mathbf{269.7 \pm 143.4}$ & 5/5 \\
Greedy-$\chi$ & $402.4 \pm 50.5$ & 5/5 \\
\bottomrule
\end{tabular}
\end{table}

\begin{figure}[t]
\centering
\includegraphics[width=\columnwidth]{figures/fig_plan_b_results.pdf}
\caption{Diagnostic and budgeting results. (a)~Equal-budget Nav ranking flip:
after the first batch, early $J$ ranks disagree with the oracle terminal
ordering (seed~2). (b)~Planar5 equal-budget cost trajectories: UP improves
fastest early but plateaus; DOWN overtakes and reaches a deeper basin.
(c)~Planar5 policy suite ($n=5$): adaptive $s_r$ policies beat Uniform;
Greedy-$\chi$ remains worse despite full feasibility.}
\label{fig:plan_b_results}
\end{figure}

Fig.~\ref{fig:plan_b_results} summarises the ranking-flip cost trajectories
and the Planar5 policy comparison underlying Tables~\ref{tab:ranking_flip}
and~\ref{tab:b3_results}.

\subsubsection{Computational Overhead}

GP inference in BO adds $\sim\!1.2$~ms per scheduling round (Cholesky of an
$8\times8$ matrix, $O(K^3)$), negligible compared to per-iteration iLQR cost
(under 1~ms for B1 and B3).  The overhead of $\chi$ computation adds no
measurable latency.  Per-run wall times are dominated by iLQR iterations:
$\sim\!10$~s for B1, $\sim\!0.3$~s for B3.

The solver is implemented in C++17 with Eigen~3.4; experiments run on an
AMD Ryzen~7~5800X (8-core, 3.8~GHz) with 32~GB RAM, Ubuntu~22.04.
\section{Discussion and Limitations}
% ===========================================================================
We have presented a feasibility-aware seed-budgeting framework for
constrained iLQR, grounded in a ranking-inversion analysis
(Corollary~\ref{cor:greedy_insufficiency}) and a protocol that treats
$\chi$ as a termination criterion rather than the primary reward.
On two independent systems---Nav2D--Hard (wheeled navigation with dense
obstacles) and Planar5 (5-DOF arm with UP/DOWN homotopies)---adaptive
allocation with $s_r$ improves over Uniform in feasible-first ranking.
The negative control Greedy-$\chi$ either loses feasibility or attains
higher cost, confirming that $\chi$ is not a reliable allocation signal.
Ranking inversion is further demonstrated on a contact-rich quadruped
template that is distinct from both the navigation and arm domains.

Several limitations warrant discussion.  \textbf{Policy separation.}
UCB, BO--EI, and Greedy-$s_r$ remain tied on all three benchmarks;
during warm-up the best seed is already identifiable by $s_r$, and the
exploration bonus does not yield additional improvement on these
libraries.  Harder seed libraries with deeper local minima---where the
greedy choice is actively misleading---would be needed to differentiate
the policies.  Similarly, the quadruped bandit suite is near-tied because
warm-up already surfaces the better basin; the quadruped is cited as a
ranking-flip diagnostic, not a policy-separation benchmark.  On
Nav2D--Hard, few runs return fully feasible incumbents at
$\varepsilon=10^{-3}$; we therefore report feasible-first comparisons
rather than a single percentage gap.  In one of the five seeds the
adaptive winner is decided by lower violation alone despite higher cost;
the feasible-first ordering is deliberately conservative about cost under
infeasibility.

\textbf{Scope of $s_r$.}
The reward $-\Delta J_{\mathrm{feas}} - \lambda\,\mathrm{viol}_+ + \eta\rho$
is a heuristic scalarisation.  Its dominant term belongs to $S_{\mathrm{LI}}$,
so ranking inversion is guaranteed at the abstract level
(Lemma~\ref{lem:li_incompleteness}).  The violation penalty lacks a
provable separation guarantee; its effectiveness depends on the constraint
landscape and is verified empirically.

\textbf{Generality.}
The framework supports arbitrary $K$, budget $B$, batch size $b$, and policy
backends; the same scheduler can host any constrained iLQR solver without
architectural changes.  Extensions to DDP and SQP are straightforward.

% ===========================================================================
\section{Conclusion and Future Work}
% ===========================================================================
We formulate multi-start constrained iLQR as \emph{feasibility-aware seed
budgeting}: allocate batches across warm-starts using diagnostics of feasible
cost improvement, violation, and rate, with $\chi$ as a termination criterion
rather than the primary reward.
Three theoretical layers---$\chi$ non-monotonicity (Section~\ref{sec:chi_analysis}),
ranking inversion analysis (Corollary~\ref{cor:greedy_insufficiency}),
and constructive inversion (Lemma~\ref{lem:li_incompleteness})---
establish that single-batch local-improvement scores are not universally
optimal guides for seed allocation.  Uniform, UCB, and GP--EI policies
share an open, budget-exact scheduler hosted on a standard constrained
iLQR stack.

On two independent systems (Nav2D--Hard and Planar5), adaptive allocation
with $s_r$ improves over Uniform under feasible-first ranking, while
Greedy-$\chi$ either loses feasibility or attains higher cost.  Ranking
inversion is confirmed on three physically distinct domains---wheeled
navigation, legged quadruped, and articulated arm---supporting the
theoretical motivation across diverse robotics platforms.
The ranking-flip evidence (Table~\ref{tab:ranking_flip}) and the
contact-rich diagnostic (Table~\ref{tab:quad_diag}) instantiate the
local-improvement incompleteness argument on real constrained-iLQR
trajectories.

Future work includes richer seed generators, optional MOBO when speed and
quality conflict, solver-agnostic hosting on alternative
solvers via the same scheduler API, and online MPC receding-horizon
re-budgeting.  The benchmark definitions and all allocation traces are
released at
\url{https://github.com/jiesu-zju/optimal-solver}; the seed-budgeting
scheduler and the constrained-iLQR host are provided as prebuilt
binaries, and every reported experiment can be reproduced end-to-end with
the exact command lines and seed parameters documented in the release.

% ===========================================================================
% ===========================================================================
\begin{thebibliography}{99}
\bibliographystyle{IEEEtran}

\bibitem{tassa2012}
Y.~Tassa, T.~Erez, and E.~Todorov, ``Synthesis and stabilization of complex behaviors through online trajectory optimization,''
in \textit{Proc. IEEE/RSJ Int. Conf. Intell. Robots Syst. (IROS)}, 2012, pp.~4906--4913.

\bibitem{howell2019}
T.~A.~Howell, B.~E.~Jackson, and Z.~Manchester, ``ALTRO: A fast solver for constrained trajectory optimization,''
in \textit{Proc. IEEE/RSJ Int. Conf. Intell. Robots Syst. (IROS)}, 2019, pp.~7674--7681.

\bibitem{mastalli2020}
C.~Mastalli et al., ``Crocoddyl: An efficient and versatile framework for multi-contact optimal control,''
in \textit{Proc. IEEE Int. Conf. Robot. Autom. (ICRA)}, 2020, pp.~2536--2542.

\bibitem{zhu2024}
Y.~Zhu, J.~J.~Payne, S.~Feng, and C.~G.~Atkeson, ``Convergent iLQR for Safe Trajectory Planning and Control of Legged Robots,''
in \textit{Proc. IEEE Int. Conf. Robot. Autom. (ICRA)}, 2024.

\bibitem{mckay1979}
M.~D.~McKay, R.~J.~Beckman, and W.~J.~Conover, ``A comparison of three methods for selecting values of input variables in the analysis of output from a computer code,''
\textit{Technometrics}, vol.~21, no.~2, pp.~239--245, 1979.

\bibitem{gyorgy2011}
A.~Gy\"{o}rgy and L.~Kocsis, ``Efficient multi-start strategies for local search algorithms,''
\textit{J. Artif. Intell. Res.}, vol.~41, pp.~407--444, 2011.

\bibitem{jain2019}
S.~Jain, D.~Salau, and A.~Agrawal, ``A review of multi-start optimization algorithms,''
\textit{Int. J. Res. Appl. Sci. Eng. Technol.}, vol.~7, no.~5, pp.~2452--2457, 2019.

\bibitem{auer2002}
P.~Auer, N.~Cesa-Bianchi, and P.~Fischer, ``Finite-time analysis of the multiarmed bandit problem,''
\textit{Mach. Learn.}, vol.~47, no.~2, pp.~235--256, 2002.

\bibitem{srinivas2010}
N.~Srinivas, A.~Krause, S.~Kakade, and M.~Seeger, ``Gaussian process optimization in the bandit setting: No regret and experimental design,''
in \textit{Proc. Int. Conf. Mach. Learn. (ICML)}, 2010, pp.~1015--1022.

\bibitem{shahriari2016}
B.~Shahriari, K.~Swersky, Z.~Wang, R.~P.~Adams, and N.~de~Freitas, ``Taking the human out of the loop: A review of Bayesian optimization,''
\textit{Proc. IEEE}, vol.~104, no.~1, pp.~148--175, 2016.

\bibitem{merkt2021}
W.~Merkt, V.~Ivan, and S.~Vijayakumar, ``Fast global motion planning via pre-computed motion primitives,''
in \textit{Proc. IEEE Int. Conf. Robot. Autom. (ICRA)}, 2021, pp.~7873--7879.

\bibitem{posa2014}
M.~Posa, C.~Cant\={u}, and R.~Tedrake, ``A direct method for trajectory optimization of rigid bodies through contact,''
\textit{Int. J. Robot. Res.}, vol.~33, no.~1, pp.~69--81, 2014.

\bibitem{schulman2014}
J.~Schulman et al., ``Motion planning with sequential convex optimization and convex collision checking,''
\textit{Int. J. Robot. Res.}, vol.~33, no.~9, pp.~1251--1270, 2014.

\bibitem{rasmussen2006}
C.~E.~Rasmussen and C.~K.~I.~Williams, \textit{Gaussian Processes for Machine Learning}.
Cambridge, MA, USA: MIT Press, 2006.

\end{thebibliography}

\end{document}
